\documentclass[11pt,a4paper]{article}
\usepackage[margin=25mm]{geometry}
\usepackage{amsmath,amssymb,amsthm,booktabs,hyperref}
\setlength{\emergencystretch}{2em}
\newtheorem{lemma}{Lemma}\newtheorem{theorem}[lemma]{Theorem}
\newcommand{\gen}[1]{\langle#1\rangle}
\newcommand{\Sym}{\operatorname{Sym}}
\newcommand{\Alt}{\operatorname{Alt}}
\title{The first critical symmetric factors}
\author{Human Compression project: a component of the L019 recovery}
\date{}
\begin{document}\maketitle
We retain the actual one-point attachment reduction. The remaining
configuration has a fractured natural symmetric factor $P$ of even
degree $m$ and rank $m/2+1$, and an even six-point rank-three $A_5$
factor $E$. Their supports meet only in the endpoint anchor $a$ and
their labels are consecutive disjoint intervals. This note settles
$m\le8$. It makes no assertion about the remaining even degrees.

\begin{lemma}\label{six}
A fractured natural $S_6$ string C-marking of rank four with an
endpoint anchor has an internal perfect split with one trivial
internal factor and with both attaching groups odd.
\end{lemma}
\begin{proof}
Reverse the marks and normalize the anchor to $1$. Under its point
stabilizer the first mark has one of the three representatives
\[
          (1\ 2),\quad (1\ 2)(3\ 4),\quad
                         (1\ 2)(3\ 4)(5\ 6).
\]
The earlier complete bounded computation tries every involution of
$(S_6)_1$ at each later position. It requires distant commutation
and adjacent noncommutation. The latter is necessary because the
anchor forces the noncommutation diagram to be connected. It checks
every pair of subset intersections, and requires the full group to
have order $720$. Thus it omits no anchored marking up to conjugacy;
it need not identify all remaining equivalent tuples.

There are $32$ terminal tuples in this enumeration, of which $24$
have every maximal omission intransitive. For each such tuple and
each internal position $j=1,2$, the verifier computes the actual
omission orbits. It tries both choices of their order $A,B$ and
requires the prefix to be transitive on $A$ and fix $B$ pointwise,
and the suffix to be transitive on $B$ and fix $A$ pointwise.
It then requires a unique crossing $a_jb_j$ of the omitted mark.
It tests pointwise that the remaining action on at least one side
is trivial. When the $A$-action is trivial it tests oddness of the
prefix and of the group generated by the crossing mark and suffix;
when the $B$-action is trivial it uses the reversed pair of groups.
Oddness is equivalent to at least one odd distinguished generator.
All $24$ fractured cases have a position passing these tests.
The unchanged complete script, terminal receipt and output are
retained from the earlier work; this proof does not extrapolate
from a selected example. The tests are exactly the asserted
conditions, proving the lemma.
\end{proof}

\begin{lemma}\label{eight}
Let $(t_0,\ldots,t_4)$ be a fractured natural $S_8$ C-marking,
with $1$ moved only by $t_0$. If it occurs in an actual one-point
attachment at $1$ to a disjointly labelled transitive factor
on at least two points, then
$t_0$ is a pure transposition.
\end{lemma}
\begin{proof}
Here the external factor need not be odd. Its original intersection
property already gives the acquired even-kernel conditions
\[
 P_j\cap\Alt(X_{j-1})\le P_{j-1},\qquad
 P_j=\gen{t_0,\ldots,t_{j-1}},\quad X_j=1^{P_j}.
                                                               \tag{1}
\]
We inspect the complete anchored degree-eight rank-five enumeration.
Its search is the earlier full even-wreath enumerator with ambient
group replaced by $S_8$ and target rank changed to five. It lists
all nonidentity involutions literally, fixes point $1$ after the
first mark, and uses the exact point stabilizer and successive
centralizers for conjugacy reduction. It tests all prefix subset
intersections and the necessary order bound
$[S_8:H]\ge2^{5-l}$ for a prefix of length $l$.
Its completeness proof is unchanged: each conjugation fixes the
previous marks and anchor, and every independent next mark at
least doubles the group order. At termination it requires order
$40320$. The node counts, including the empty prefix, are
\[
                    1,4,60,432,90,22.
\]
Exactly nine terminal tuples generate $S_8$. The first six have
a pure first mark. The ninth has a transitive maximal omission,
so is excluded by fracture. The seventh and eighth have, with
one-based table numbering, the following exact supported kernels:
\begin{center}
\begin{tabular}{crrrr}\toprule
Tuple&$j$&$|P_j\cap\Sym(X_{j-1})|$
 &$|P_{j-1}\cap\Sym(X_{j-1})|$
 &$|P_j\cap\Alt(X_{j-1})|$\\\midrule
7&3&120&10&60\\
8&4&720&60&360\\\bottomrule
\end{tabular}
\end{center}
Their previous supported kernels consist entirely of even permutations.
Thus both violate (1), by index six even before their odd elements
are considered. A separate literal-permutation implementation
enumerates the groups, omission orbits and supported kernels of all
nine tuples and verifies these orders and inclusions. This excludes
the only two nonpure fractured candidates and proves the lemma.
\end{proof}

\begin{theorem}
In the actual remaining $P,E$ configuration described above, if
$m\le8$, the whole marking has an internal sufficient split and hence
a valid rank and degree extension. For $m=6$ the sufficient split is
inherited from $P$; for $m=8$ the surviving pure first mark gives it.
The only critical symmetric degrees still unpaid are $m\ge10$.
\end{theorem}
\begin{proof}
For $m=6$, Lemma~\ref{six} and the acquired one-point transport
lemma give a sufficient internal perfect split in the whole marking.
The acquired odd-attachment theorem then supplies its extension.

For $m=8$, Lemma~\ref{eight} applies with the actual external $E$.
Write the pure first mark of $P$ as $p=(a\ c)$.
The tail of $P$ fixes $a$ and is transitive on $X\setminus\{a\}$:
otherwise adjoining the sole edge $ac$ cannot make $P$ transitive.
In the whole marking, omission of $p$ therefore has exactly the
orbits $Y$ and $X\setminus\{a\}$, with the two corresponding
label intervals supported and transitive there. Thus $p$ itself is
an internal pure perfect split, which is sufficient in the precise
sense used by the acquired induction.

For completeness, $m=2$ would require two independent marks in
$S_2$ and is impossible. At $m=4$, normalize the anchor to $1$.
The last two marks are transpositions on the other three points
and must not commute. Write them $(b\ c),(c\ d)$ after relabelling.
The first mark moves $1$ and commutes with $(c\ d)$; it is either
$(1\ b)$ or $(1\ b)(c\ d)$. In the second case, the first two
marks act transitively on the four points, contradicting fracture
of the last omission. In the first case all three marks form the
simplex path and its middle mark is an internal pure perfect split.
Transport and the acquired pure-split theorem finish this case.
These are all positive even $m<10$.
\end{proof}

\paragraph{Attribution and boundary.}
The principal source is Cameron--Fernandes--Leemans,
\emph{The number of string C-groups of high rank},
arXiv:2212.12723v2. We acquire the reviewed one-point transport,
even-kernel and exact leaf criteria, and the six-point $A_5$
reduction; their original group-theoretic and software premises
remain bound. The degree-six search is reused unchanged. An earlier
degree-eight centralizer diagnostic traversed this same rank but
retained only centralizer counterexamples; the present complete
terminal list and literal leaf filter are retained separately.
No conjecture, whole L019, Lean theorem or complexity measurement
is claimed. The general critical leaf argument and the reversible
stability correspondence are still required.
\end{document}
